\section{Vulnerability Replay Evidence and Rediscovery Protocol}
\label{sec:cve}
The vulnerability study comprises four cases: three associated with CVE identifiers and one SMARTPL availability failure without a CVE. Each case has a report and a triggering input; three have direct execution evidence, and ProFTPD also has a reported patched-version comparison. Tables~\ref{tab:cve} and~\ref{tab:cveresults} summarize these observations. Case identities, versions, and replay counts follow the original reports; incomplete patch and experimental-configuration evidence prevents independent verification of every version pair.

\subsection{Case identity and evidence provenance}
We distinguish a direct execution observation (L) from a replay result stated only in a case report (R). A report can contain additional trials not represented by the available execution evidence. Duplicate diagnostic observations do not constitute independent trials. Input identity and byte length are retained in the supporting evidence. This retrospective analysis involves no additional target executions. Table~\ref{tab:cve} identifies the reported versions, triggering conditions, and comparative evidence.

\begin{table*}[!tbp]
\centering\small
\caption{Vulnerability cases and supporting evidence (L: direct observation; R: report).}
\label{tab:cve}
\begin{tabularx}{\textwidth}{@{}>{\raggedright\arraybackslash}p{0.20\textwidth} >{\raggedright\arraybackslash}p{0.19\textwidth} Y Y@{}}
\toprule
Case / interface & Reported vulnerable version & Trigger evidence / oracle & Patch or control evidence \\
\midrule
\mbox{CVE-2023-51713}\newline ProFTPD / FTP & Reported vulnerable version; ASan; 65,535-byte command buffer & 60,006-byte input; heap-buffer-overflow: one-byte read beyond a 65,568-byte region (L). & Reported security patch; three rejections, no ASan (L); same input/buffer. \\
\addlinespace[3pt]
\mbox{CVE-2026-38998}\newline LIVE555 / RTSP & Campaign: 2023.05.10; public report: 2026.02.26 (R) & 728-byte public PoC; heap-use-after-free (L). Fuzzer-discovered input absent. & Attempted local patch: 1/1 still triggers (R); upstream-patch comparison absent. \\
\addlinespace[3pt]
\mbox{CVE-2026-39863}\newline Kamailio / SIP over TCP & Reported vulnerable version; TCP required & 278-byte input; parsed Content-Length $-10$ (L): overflow predicate, not crash count. & Local overflow guard suppresses diagnostic (R); equivalence to the identified security patch unverified. \\
\addlinespace[3pt]
SMARTPL availability case\newline forked-daapd / HTTP & Reported version 27.2 & 148-byte input; lexer errors; health probe fails; process alive (R). Direct replay evidence absent. & No patched replay/CVE; parser reportedly removed in 28.4. \\
\bottomrule
\end{tabularx}
\end{table*}

For ProFTPD, the reported defect is an out-of-bounds read during FTP command processing, supported by an ASan heap-buffer-overflow observation. The oracle uses a 65,535-byte command buffer; the report states that the default 512-byte configuration does not expose the same error to ASan. The result is therefore conditional on this setting. Three reported patched-version trials return an FTP 500 response without ASan activation and terminate normally. This supports differential replay evidence, although experimental-version identity and exact patch equivalence remain unverified.

For LIVE555, direct diagnostic evidence supports a use-after-free triggered by the public PoC. A duplicated observation is treated as one diagnostic signature rather than independent replication. The report attributes fuzzer-generated inputs to the same memory-lifetime failure, records 8/9 replay hits, and describes an unsuccessful local patch. The fuzzer-generated reproducer, complete memory-lifetime evidence, and individual trial outcomes are unavailable. The public PoC supports replay validation, but does not demonstrate independent rediscovery by the evaluated controller or independently establish cross-version CVE attribution.

For Kamailio, a negative parsed Content-Length supports the reported integer-overflow condition. The trigger requires TCP and is not assessed by the UDP-only benchmark condition. A local overflow guard reportedly suppresses the diagnostic, but direct patched-version evidence and repetition counts are unavailable. The observation remains report-only and does not establish equivalence to the upstream patch. An earlier unsupported 20/20 versus 0/20 claim is excluded from quantitative results. The SMARTPL case reports reproducible unavailability from a known input, while none of five sampled campaign hangs satisfies that condition; neither a CVE identity nor a verified patched counterpart is established.

\subsection{Replay outcomes and their denominators}
Table~\ref{tab:cveresults} preserves each reported batch separately. A numerator counts replay attempts or selected queue inputs satisfying that case's predicate; the denominator is the number tested. These are not independent 12-hour fuzzer runs, and the main experiment's nominal $N=10$ is not substituted for the actual replay denominator. Selected batches, per-input repeats, and later aggregate counts may overlap; they are not pooled into one success rate. In particular, 19/20 selected inputs is not 95\% campaign recall, and 8/9 LIVE555 replays is below the prospective 95\% stability threshold.

\begin{table*}[!tbp]
\centering\small
\caption{Reported replay batches and separate controls; original denominators retained.}
\label{tab:cveresults}
\begin{tabularx}{\textwidth}{@{}>{\raggedright\arraybackslash}p{0.20\textwidth} Y Y Y@{}}
\toprule
Case & Vulnerable-side batches (R) & Patched or diagnostic control & Interpretation / missing evidence \\
\midrule
\mbox{CVE-2023-51713} & 1/1 seed; 1/1 fuzzer crash (Sept.~22); 3/3 batch crashes (Sept.~25). & Reported fix: 0/3 ASan activations; 3/3 FTP 500 rejections (L, Sept.~28). & Separate replay batches; no arm-specific discovery probability or trigger time. \\
\addlinespace[3pt]
\mbox{CVE-2026-38998} & 1/1 public PoC; three fuzzer inputs: 3/3 each; later aggregate: 8/9. & Local patch: 1/1 still triggers (R); successful patched-side denominator absent. & Unpooled batches; fuzzer input/trial records absent; no confirmed patch pair. \\
\addlinespace[3pt]
\mbox{CVE-2026-39863} & 1/1 seed; queue samples: 3/4 (60 s), 19/20 (60 min), 19/20 (1,540 min), 50/51 (separate 60-min batch). & Local guard: fingerprint absent (R); repeat count unavailable. & Queue samples, not campaign recall; no extrapolation to 195 entries; patched-version evidence absent. \\
\addlinespace[3pt]
SMARTPL availability case & 5/5 advisory repeats; 3/3 48-message-seed repeats; 1/1 later advisory replay. & Sampled campaign hangs: 0/5 SMARTPL hits (R); patched replay absent. & Known-input unavailability; no confirmed campaign rediscovery or CVE. \\
\bottomrule
\end{tabularx}
\end{table*}

Additional reported crashes and campaigns cannot be linked to the primary benchmark and ablation observations. Crash totals do not count unique vulnerabilities, and durations such as 60 or 1,540 minutes do not specify time to first trigger. Associations among reproducing inputs, verified A--E conditions, model information exposure, and independent root causes are incomplete. The replay outcomes therefore do not establish arm-specific discovery probabilities, novel CVEs, or admission/calibration gains.

\subsection{Controlled rediscovery still to be measured}
Confirmatory rediscovery requires a vulnerable version and the same version with a minimal security patch, holding dependencies, instrumentation, settings, and input compatibility constant. Each case needs a network-reachable trigger, independent reached/triggered criteria, reproducible initial conditions, and a validated patch comparison. The proposed six-case pilot across at least three protocols, followed by expansion toward twelve, remains incomplete. Eligibility requires at least 95\% observed oracle repeatability in a declared validation set, no activation after patching, reproducible experimental preparation, and information separation. The four cases do not constitute six qualified version pairs; a small all-success replay batch does not establish an underlying success probability of at least 95\%.

Information-controlled campaigns must exclude CVE identifiers, advisories, public PoCs, and patch contents from model inputs and campaign seeds, while retaining anonymous oracle IDs in separate benchmark metadata. The public PoC and known triggers in this collection are validation inputs; their replay is not evidence of information-controlled rediscovery. Source-level trigger conditions, sanitizer reports with patch validation, and protocol invariants distinguish reaching relevant code from actually triggering a defect. Experiments require isolated target networks and reset between trials; these are protocol requirements, not properties inferred from the reports.

For each CVE and arm, the formal plan retains ten independent runs and a 12-hour deadline, prespecified before execution. Deadline recall, time-to-trigger, reached-but-not-triggered counts, patched-target activation, executions/coverage/tokens before trigger, and recall per CPU-hour or 100,000 tokens remain unavailable for the controlled A--E comparison. Nontriggers require right-censoring records, and interrupted runs require separate failure accounting. The available replay batches do not define those risk sets, so Kaplan--Meier estimates or arm-level security-effectiveness claims are not reported.
